\exercisegroup

\textbf{1.} Prove the formula

\[
\sum_ {k = q + 1} ^ {n - p + q + 1} \binom {n - k} {p - q - 1} \binom {k - 1} {q} = \binom {n} {p},
\]

where \(p \leqslant n\) and \(q < p\) (generalize the argument of no. 8, Corollary to Proposition 14).

\begin{enumerate}
\def\labelenumi{\arabic{enumi}.}
\setcounter{enumi}{1}
\tightlist
\item
  If \(n \geqslant 1\) , prove the relation
\end{enumerate}

\[
\binom {n} {0} - \binom {n} {1} + \binom {n} {2} - \dots + (- 1) ^ {n} \binom {n} {n} = 0.
\]

(Define a one-to-one correspondence between the set of subsets of \([1, n]\) which have an even number of elements, and the set of subsets of \([1, n]\) which have an odd number of elements. Distinguish between the cases n even and n odd.)

\textbf{3.} Prove the relations

\[
\begin{array}{l} \binom {n} {0} \binom {n} {p} + \binom {n} {1} \binom {n - 1} {p - 1} + \binom {n} {2} \binom {n - 2} {p - 2} + \dots + \binom {n} {p} \binom {n - p} {0} = 2 ^ {p} \binom {n} {p}, \\ \binom {n} {0} \binom {n} {p} - \binom {n} {1} \binom {n - 1} {p - 1} + \binom {n} {2} \binom {n - 2} {p - 2} - \dots + (- 1) ^ {p} \binom {n} {p} \binom {n - p} {0} = 0. \end{array}
\]

(Consider the subsets of \(p\) elements of {[}1, n{]} which contain a given subset of \(k\) elements ( \(0 \leqslant k \leqslant p\) ), and use Exercise 2 for the second formula.)

\begin{enumerate}
\def\labelenumi{\arabic{enumi}.}
\setcounter{enumi}{3}
\tightlist
\item
  Prove Proposition 15 of no. 8 by defining a bijection of the set of mappings \(u\) of \([1, h]\) into \([0, n]\) such that
\end{enumerate}

\[
\sum_ {i = 1} ^ {h} u (x) \leqslant n
\]

onto the set of strictly increasing mappings of \([1, h]\) into \([1, n + h]\) .

\begin{enumerate}
\def\labelenumi{\arabic{enumi}.}
\setcounter{enumi}{4}
\tightlist
\item
  * (a) Let E be a distributive lattice and f be a mapping of E into a commutative semigroup M (written additively) such that
\end{enumerate}

\[
f (x) + f (y) = f (\sup (x, y)) + f (\inf (x, y))
\]

for all \(x, y\) in \(\mathbf{E}\) . Show that for each finite subset \(\mathbf{I}\) of \(\mathbf{E}\) we have

\[
\begin{array}{r l} f (\sup (I)) + \sum_ {2 n \leqslant \operatorname{Card} (I)} & \left(\sum_ {H \subset I, \operatorname{Card} (H) = 2 n} f (\inf (H))\right) \\ & = \sum_ {2 n + 1 \leqslant \operatorname{Card} (I)} \left(\sum_ {H \subset I, \operatorname{Card} (H) = 2 n + 1} f (\inf (H))\right) \end{array}
\]

(By induction on Card (I).) *

\begin{enumerate}
\def\labelenumi{(\alph{enumi})}
\setcounter{enumi}{1}
\tightlist
\item
  In particular, let A be a set, let \((\mathbf{B}_i)_{i\in \mathbf{I}}\) be a finite family of finite subsets of A, and let B be the union of the \(\mathbf{B}_i\) . For each subset H of I, put \(\mathbf{B}_{\mathbf{H}} = \bigcap_{i\in \mathbf{H}}\mathbf{B}_i\) . Show that
\end{enumerate}

\[
\begin{array}{r l} \operatorname{Card} (\mathrm{B}) + \sum_ {2 n \leqslant \operatorname{Card} (\mathrm{I})} & \left(\sum_ {\operatorname{Card} (\mathrm{H}) = 2 n} \operatorname{Card} \left(\mathrm{B} _ {\mathrm{H}}\right)\right) \\ & = \sum_ {2 n + 1 \leqslant \operatorname{Card} (\mathrm{I})} \left(\sum_ {\operatorname{Card} (\mathrm{H}) = 2 n + 1} \operatorname{Card} \left(\mathrm{B} _ {\mathrm{H}}\right)\right). \end{array}
\]

\begin{enumerate}
\def\labelenumi{\arabic{enumi}.}
\setcounter{enumi}{5}
\tightlist
\item
  Prove the formula
\end{enumerate}

\[
\begin{array}{r l} \binom {n + h} {h} = 1 + \binom {h} {1} \binom {n + h - 1} {h} - \binom {h} {2} \binom {n + h - 2} {h} + \dots \\ & + (- 1) ^ {h} \binom {h} {h} \binom {n} {h}. \end{array}
\]

(If \(\mathbf{F}\) denotes the set of mappings \(u\) of \([1, h]\) into \([0, n]\) such that

\[
\sum_ {x = 1} ^ {h} u (x) \leqslant n,
\]

consider for each subset H of {[}1, h{]} the set of all \(u \in F\) such that \(u(x) \geqslant 1\) for each \(x \in H\) , and use Exercise 5.)

\begin{enumerate}
\def\labelenumi{\arabic{enumi}.}
\setcounter{enumi}{6}
\tightlist
\item
  \begin{enumerate}
  \def\labelenumii{(\alph{enumii})}
  \tightlist
  \item
    Let \(S_{n,p}\) denote the number of mappings of \([1, n]\) onto \([1, p]\) . Prove that
  \end{enumerate}
\end{enumerate}

\[
\mathrm{S} _ {n, p} = p ^ {n} - \binom {p} {1} (p - 1) ^ {n} + \binom {p} {2} (p - 2) ^ {n} - \dots + (- 1) ^ {p - 1} \binom {p} {p - 1}.
\]

( Note that \(p^{n} = S_{n,p} + \binom{p}{1} S_{n,p-1} + \binom{p}{2} S_{n,p-2} + \cdots + \binom{p}{p-1}\) and use Exercise 3.)

\begin{enumerate}
\def\labelenumi{(\alph{enumi})}
\setcounter{enumi}{1}
\item
  Prove that \(S_{n,p} = p(S_{n-1,p} + S_{n-1,p-1})\) (method of no. 8, Proposition 13).
\item
  Prove that
\end{enumerate}

\[
\mathrm{S} _ {n + 1, n} = \frac {n}{2} (n + 1)! \quad \text { and } \quad \mathrm{S} _ {n + 2, n} = \frac {n (3 n + 1)}{2 4} (n + 2)!
\]

(consider the elements \(r\) of \([1, n]\) whose inverse image consists of more than one element).

\begin{enumerate}
\def\labelenumi{(\alph{enumi})}
\setcounter{enumi}{3}
\tightlist
\item
  If \(P_{n,p}\) is the number of partitions into p parts of a set of n elements, show that \(S_{n,p} = p!P_{n,p}\) .
\end{enumerate}

\begin{enumerate}
\def\labelenumi{\arabic{enumi}.}
\setcounter{enumi}{7}
\tightlist
\item
  Let \(p_n\) be the number of permutations \(u\) of a set \(\mathbf{E}\) with \(n\) elements such that \(u(x) \neq x\) for all \(x \in \mathbf{E}\) . Show that
\end{enumerate}

\[
p _ {n} = n! - \binom {n} {1} (n - 1)! + \binom {n} {2} (n - 2)! - \dots + (- 1) ^ {n}
\]

* and hence that \(p_n \sim n! / e\) as \(n \to +\infty\) (same method as in Exercise 7 (a)).

\begin{enumerate}
\def\labelenumi{\arabic{enumi}.}
\setcounter{enumi}{8}
\tightlist
\item
  \begin{enumerate}
  \def\labelenumii{(\alph{enumii})}
  \tightlist
  \item
    Let E be a set with qn elements. Show that the number of partitions of E into n subsets each of q elements is equal to
  \end{enumerate}
\end{enumerate}

\[
(q n)! / (n! (q!) ^ {n}).
\]

\begin{enumerate}
\def\labelenumi{(\alph{enumi})}
\setcounter{enumi}{1}
\tightlist
\item
  Suppose that \(\mathbf{E} = [1, qn]\) . Show that the number of partitions of \(\mathbf{E}\) into \(n\) subsets of \(q\) elements, no one of which is an interval, is equal to
\end{enumerate}

\[
\frac {(q n) !}{n ! (q !) ^ {n}} - \frac {(q n - q + 1) !}{1 ! (n - 1) ! (q !) ^ {n - 1}} + \frac {(q n - 2 q + 2) !}{2 ! (n - 2) ! (q !) ^ {n - 2}} - \dots + (- 1) ^ {n}
\]

(same method as in Exercises 7 and 8).

\begin{enumerate}
\def\labelenumi{\arabic{enumi}.}
\setcounter{enumi}{9}
\tightlist
\item
  Let \(q_{n,k}\) be the number of strictly increasing mappings \(u\) of \([1, k]\) into \([1, n]\) such that for each even (resp. odd) \(x\) , \(u(x)\) is even (resp. odd). Show that \(q_{n,k} = q_{n-1, k-1} + q_{n-2, k}\) and deduce that
\end{enumerate}

\[
q _ {n, k} = \binom {\left[ \frac {n + k}{2} \right]} {k}.
\]

\begin{enumerate}
\def\labelenumi{\arabic{enumi}.}
\setcounter{enumi}{10}
\tightlist
\item
  Let E be a set with n elements and let S be a set of signs such that S is the disjoint union of E and a set consisting of a single element f.~Suppose that f has weight 2 and that each element of E has weight 0 (Chapter I, Appendix, Exercise 3).
\end{enumerate}

\begin{enumerate}
\def\labelenumi{(\alph{enumi})}
\tightlist
\item
  Let M be the set of significant words in \(L_{0}(S)\) which contain each element of E exactly once. Show that if \(u_{n}\) is the number of elements in M, then \(u_{n+1} = (4n - 2)u_{n}\) , and deduce that
\end{enumerate}

\[
u _ {n} = 2. 6 \dots (4 n - 6) \quad (n \geqslant 2).
\]

(This is the number of products of n different terms with respect to a non-associative law of composition.)

\begin{enumerate}
\def\labelenumi{(\alph{enumi})}
\setcounter{enumi}{1}
\tightlist
\item
  Let \(x_{i}\) be the \(i\) th of the elements of \(E\) which appears in a word of \(M\) . Show that the number \(v_{n}\) of words of \(M\) , for which the sequence \((x_{i})\) is given, is equal to \((\frac{2n-2}{n-1})/n\) and satisfies the relation
\end{enumerate}

\[
v _ {n + 1} = v _ {1} v _ {n} + v _ {2} v _ {n - 1} + \dots + v _ {n - 1} v _ {2} + v _ {n} v _ {1}.
\]

\begin{enumerate}
\def\labelenumi{\arabic{enumi}.}
\setcounter{enumi}{11}
\tightlist
\item
  \begin{enumerate}
  \def\labelenumii{(\alph{enumii})}
  \tightlist
  \item
    Let \(p\) and \(q\) be two integers \(\geqslant 1\) , let \(n = 2p + q\) , let \(E\) be a set with \(n\) elements, and let \(N = \binom{n}{p} = \binom{n}{p + q}\) . Let \((X_i)_{1 \leqslant i \leqslant N}\) (resp. \((Y_i)_{1 \leqslant i \leqslant N}\) ) be the sequence of all subsets of \(E\) which have \(p\) (resp. \(p + q\) ) elements arranged in a certain order. Show that there exists a bijection \(\varphi\) of {[}1, N{]} onto itself such that \(X_{\varphi(i)} \subset Y_i\) for all \(i\) . (The method is analogous to that of Exercise 6 of §4: observe that for each \(r \leqslant N\) the number of sets \(Y_j\) which contain at least one of \(X_1, \ldots, X_r\) is \(\geqslant r\) .)
  \end{enumerate}
\end{enumerate}

\begin{enumerate}
\def\labelenumi{(\alph{enumi})}
\setcounter{enumi}{1}
\tightlist
\item
  Let \(h, k\) be two integers \(\geqslant 1\) , let \(n\) be an integer such that \(2h + k < n\) , let \(\mathbf{E}\) be a set with \(n\) elements and let \((\mathbf{X}_i)_{1 \leqslant i \leqslant r}\) be a sequence of distinct subsets of \(\mathbf{E}\) , each having \(h\) elements. Show that there exists a sequence \((\mathbf{Y}_j)_{1 \leqslant j \leqslant r+1}\) of distinct subsets of \(\mathbf{E}\) , each having \(h + k\) elements, such that each \(\mathbf{Y}_j\) contains at least one \(\mathbf{X}_i\) and each \(\mathbf{X}_i\) is contained in at least one \(\mathbf{Y}_j\) (by induction on \(n\) , using (a)).
\end{enumerate}

¶ 13. Let E be a set with 2m elements, let q be an integer \textless{} m, and let F be the set of all subsets S of P(E) with the following property: if X and Y are two distinct elements of S such that X ⊂ Y, then Y --- X has at most 2q elements.

\begin{enumerate}
\def\labelenumi{(\alph{enumi})}
\item
  Let \(\mathfrak{M} = (\mathrm{A}_i)_{1\leqslant i\leqslant p}\) be an element of \(\mathcal{F}\) such that \(p = \operatorname{Card}(\mathfrak{M})\) is as large as possible. Show that \(m - q\leqslant \operatorname{Card}(\mathrm{A}_i)\leqslant m + q\) for \(1\leqslant i\leqslant p\) . (Argue by contradiction. Suppose, for example, that there exist indices \(i\) such that \(\operatorname{Card}(\mathrm{A}_i) < m - q\) , and consider those of the \(\mathrm{A}_i\) for which \(\operatorname{Card}(\mathrm{A}_i)\) has the least possible value \(m - q - s\) (where \(s\geqslant 1\) ). Let \(\mathrm{A}_1,\ldots ,\mathrm{A}_r\) , say, be these sets. Let \(\mathfrak{G}\) be the set of subsets of E each of which is the union of some \(\mathrm{A}_i(1\leqslant i\leqslant r)\) and a subset of \(2q + 1\) elements contained in \(\mathrm{E} - \mathrm{A}_i\) . Show that \(\mathfrak{G}\) contains at least \(r + 1\) elements (cf.~Exercise 12), and that if \(\mathrm{B}_1,\ldots ,\mathrm{B}_{r + 1}\) are \(r + 1\) distinct elements of \(\mathfrak{G}\) , the set whose elements are \(\mathrm{B}_j(1\leqslant j\leqslant r + 1)\) and \(\mathrm{A}_i(r + 1\leqslant i\leqslant p)\) belongs to \(\mathcal{F}\) , contrary to hypothesis.)
\item
  Deduce from (a) that the number of elements \(p\) of each \(\mathfrak{S} \in \mathcal{F}\) satisfies the inequality
\end{enumerate}

\[
p \leqslant \sum_ {k = 0} ^ {2 q} \binom {2 m} {m - q + k}.
\]

\begin{enumerate}
\def\labelenumi{(\alph{enumi})}
\setcounter{enumi}{2}
\tightlist
\item
  Establish results analogous to those of (a) and (b) when 2m or 2q is replaced by an uneven number.
\end{enumerate}

\begin{enumerate}
\def\labelenumi{\arabic{enumi}.}
\setcounter{enumi}{13}
\tightlist
\item
  Let E be a finite set with n elements, let \((a_{j})_{1 \leqslant j \leqslant n}\) be the sequence of elements of E arranged in some order, and let \((\mathbf{A}_{i})_{1 \leqslant i \leqslant m}\) be a sequence of subsets of E.
\end{enumerate}

\begin{enumerate}
\def\labelenumi{(\alph{enumi})}
\tightlist
\item
  For each index \(j\) , let \(k_{j}\) be the number of indices \(i\) such that \(a_{j} \in \mathbf{A}_{i}\) , and let \(\mathbf{S}_{i} = \operatorname{Card}(\mathbf{A}_{i})\) . Show that
\end{enumerate}

\[
\sum_ {j = 1} ^ {n} k _ {j} = \sum_ {i = 1} ^ {m} s _ {i}.
\]

\begin{enumerate}
\def\labelenumi{(\alph{enumi})}
\setcounter{enumi}{1}
\item
  Suppose that, for each subset \(\{x, y\}\) of two elements of E, there exists exactly one index \(i\) such that \(x\) and \(y\) are contained in \(\mathbf{A}_i\) . Show that, if \(a_j \notin \mathbf{A}_i\) , then \(\mathbf{S}_i \leqslant k_j\) .
\item
  With the hypotheses of (b), show that \(m \geqslant n\) . (Let \(k_n\) be the least of the numbers \(k_j\) . Show that we may suppose that, whenever \(i \leqslant k_n\) , \(j \leqslant k_n\) , and \(i \neq j\) , we have \(a_j \notin A_i\) and \(a_n \notin A_j\) for all \(j \geqslant k_n\) .)
\item
  With the hypotheses of (b), show that \(m = n\) if and only if one of the following two alternatives is true: (i) \(A_1 = \{a_1, a_2, \ldots, a_{n-1}\}\) , \(A_i = \{a_{i-1}, a_n\}\) for \(i = 2, \ldots, n\) ; (ii) \(n = k(k - 1) + 1\) , each \(A_i\) has \(k\) elements, and each element of \(E\) belongs to exactly \(k\) sets \(A_i\) .
\end{enumerate}

\begin{enumerate}
\def\labelenumi{\arabic{enumi}.}
\setcounter{enumi}{14}
\tightlist
\item
  Let E be a finite set, let \(\mathfrak{L}\) and \(\mathfrak{G}\) be two disjoint non-empty subsets of \(\mathfrak{P}(\mathbf{E})\) , and let \(\lambda, h, k, l\) be four integers \(\geqslant 1\) with the following properties: (i) for each \(\mathbf{A} \in \mathfrak{L}\) and each \(\mathbf{B} \in \mathfrak{G}\) , Card \((\mathbf{A} \cap \mathbf{B}) \geqslant \lambda\) ; (ii) for each \(\mathbf{A} \in \mathfrak{L}\) , Card \((\mathbf{A}) \geqslant h\) ; (iii) for each \(\mathbf{B} \in \mathfrak{G}\) , Card \((\mathbf{B}) \leqslant k\) ; (iv) for each \(x \in \mathbf{E}\) , the number of elements of \(\mathfrak{L} \cup \mathfrak{G}\) which contain \(x\) is exactly \(l\) . Show that Card \((\mathbf{E}) \leqslant hk / \lambda\) . (Let \((a_i)_{1 \leqslant i \leqslant n}\) be the sequence of distinct elements of E arranged in some order, and for each \(i\) let \(r_i\) be the number of elements of \(\mathfrak{L}\) to which \(a_i\) belongs. Show that, if Card \((\mathfrak{L}) = s\) and Card \((\mathfrak{G}) = t\) , then we have
\end{enumerate}

\[
\sum_ {i = 1} ^ {n} r _ {i} \leqslant s h, \quad \sum_ {i = 1} ^ {n} (l - r _ {i}) \leqslant t k, \quad \text { and } \quad \sum_ {i = 1} ^ {n} r _ {i} (l - r _ {i}) \geqslant \lambda s t.)
\]

For Card (E) to be equal to \(hk / \lambda\) it is necessary and sufficient that for each \(A \in \mathfrak{L}\) and each \(B \in \mathfrak{C}\) we have Card (A) = \(h\) , Card (B) = \(k\) ,

Card \((\mathbf{A} \cap \mathbf{B}) = \lambda\) , and that there exists an \(r \leqslant l\) such that for each \(x \in \mathbf{E}\) the number of elements of \(\mathfrak{L}\) to which \(x\) belongs is equal to \(r\) .

\begin{enumerate}
\def\labelenumi{\arabic{enumi}.}
\setcounter{enumi}{15}
\tightlist
\item
  Let E be a finite set with n elements, let D be a non-empty subset of \(\mathfrak{P}(\mathrm{E})\) , and let \(\lambda\) , k, l be three integers \(\geqslant 1\) with the following properties: (i) if A, B are distinct elements of D, then Card \((\mathrm{A} \cap \mathrm{B}) = \lambda\) ; (ii) for each \(A \in D\) , Card \((\mathrm{A}) \leqslant k\) ; (iii) for each \(x \in E\) the number of elements of D to which x belongs is equal to l. Show that
\end{enumerate}

\[
n (\lambda - 1) \leqslant k (k - 1),
\]

and that if \(n(\lambda - 1) = k(k - 1)\) then \(\lambda = k\) and \(\operatorname{Card}(\mathfrak{D}) = n\) . (Given \(a \in E\) , let \(\mathfrak{L}\) be the set of all \(A - \{a\}\) where \(A \in \mathfrak{D}\) and \(a \in A\) , and let \(\mathfrak{C}\) be the set of all \(A \in \mathfrak{D}\) such that \(a \notin A\) . Apply the results of Exercise 15 to \(\mathfrak{L}\) and \(\mathfrak{C}\) .)

\begin{enumerate}
\def\labelenumi{\arabic{enumi}.}
\setcounter{enumi}{16}
\tightlist
\item
  Let \(i, h, k\) be three integers such that \(i \geqslant 1\) , \(h \geqslant i\) , \(k \geqslant i\) . Show that there exists an integer \(m_i(h, k)\) with the following properties: for each finite set E with at least \(m_i(h, k)\) elements, and each partition \((\mathcal{X}, \mathcal{Y})\) of the set \(\mathfrak{F}_i(\mathbf{E})\) of subsets of \(i\) elements of E, it is impossible that every subset of \(h\) elements of E contains a subset \(\mathbf{X} \in \mathcal{X}\) and that every subset of \(k\) elements of E contains a subset \(\mathbf{Y} \in \mathcal{Y}\) ; in other words, if every subset of \(h\) elements of E contains some \(\mathbf{X} \in \mathcal{X}\) , there exists a subset A of \(k\) elements of E such that every subset of \(i\) elements of A belongs to \(\mathcal{X}\) . (Proof by induction. Show that we may take
\end{enumerate}

\[
m _ {1} (h, k) = h + k - 1, \quad m _ {i} (i, k) = k, \quad \text { and } \quad m _ {i} (h, i) = h,
\]

and finally \(m_i(h, k) = m_{i-1}(m_i(h-1, k), m_i(h, k-1)) + 1\) . If \(E\) is a set with \(m_i(h, k)\) elements, if \(a \in E\) and if \(E' = E - \{a\}\) , show that if the proposition were false, then every subset of \(m_i(h-1, k)\) elements of \(E'\) would contain a subset \(X'\) of \(i-1\) elements such that \(X' \cup \{a\} \in \mathcal{X}\) , and that every subset of \(m_i(h, k-1)\) elements of \(E'\) would contain a subset \(Y'\) of \(i-1\) elements such that \(Y' \cup \{a\} \in \mathcal{Y}\) .

\begin{enumerate}
\def\labelenumi{\arabic{enumi}.}
\setcounter{enumi}{17}
\tightlist
\item
  \begin{enumerate}
  \def\labelenumii{(\alph{enumii})}
  \tightlist
  \item
    Let E be a finite ordered set with p elements. If m, n are two integers such that mn \textless{} p, show that E has either a totally ordered subset of m elements or else a free subset (§ 1, Exercise 5) of n elements (use § 4, Exercise 5).
  \end{enumerate}
\end{enumerate}

\begin{enumerate}
\def\labelenumi{(\alph{enumi})}
\setcounter{enumi}{1}
\tightlist
\item
  Let \(h, k\) be two integers \(\geqslant 1\) and let \(r(h, k) = (h - 1)(k - 1) + 1\) . Let \(I\) be a finite totally ordered set with at least \(r(h, k)\) elements. Show that, for each finite sequence \((x_i)_{i \in I}\) of elements of a totally ordered set \(E\) , there exists either a subset \(H\) of \(h\) elements of \(I\) such that the sequence \((x_i)_{i \in H}\) is increasing, or else a subset \(K\) of \(k\) elements of \(I\) such that the sequence \((x_i)_{i \in K}\) is decreasing. (Use (a) applied to \(I \times E\) .)
\end{enumerate}

